\documentclass[11pt]{article}
\usepackage[margin=0.9in]{geometry}
\usepackage{amsmath,amssymb}
\usepackage{enumitem}
\usepackage{tcolorbox}
\usepackage{xcolor}
\usepackage{booktabs}
\usepackage{array}
\usepackage{fancyhdr}
\tcbuselibrary{skins,breakable}

\definecolor{boxbg}{gray}{0.94}
\definecolor{boxline}{gray}{0.55}
\definecolor{keybg}{gray}{0.97}

\pagestyle{fancy}
\fancyhf{}
\renewcommand{\headrulewidth}{0pt}
\fancyfoot[C]{\small\color{gray} Intro Finance \quad$\cdot$\quad The Cookie Box Pricing Decision \quad$\cdot$\quad Case Study \#2}
\fancyfoot[R]{\small\color{gray}\thepage}

\newtcolorbox{databox}{enhanced, breakable, colback=boxbg, colframe=boxline,
  boxrule=0.6pt, arc=3pt, left=10pt, right=10pt, top=6pt, bottom=6pt}

\newcommand{\question}[1]{\medskip\noindent\textbf{#1}\par\smallskip}
\newcommand{\hint}[1]{{\small\color{gray}\textit{Hint: #1}}\par}

\begin{document}
\thispagestyle{fancy}
{\small NAME: \underline{\hspace{2.6in}}}\par\vspace{4pt}
\begin{center}
{\LARGE\bfseries The Cookie Box Pricing Decision}\\[3pt]
\rule{2.2in}{0.8pt}\\[4pt]
{\itshape a small-business pricing decision \quad$\bullet$\quad Intro Finance}
\end{center}

\vspace{2pt}

Amara owns a small bakery and wants to launch a line of gourmet cookie boxes. Her accountant is blunt: the new line only makes sense if it earns a \textbf{40\% profit margin}, and it must carry \$1{,}200 a month of extra kitchen and lease costs on its own. Her business partner has already printed a price sheet using a rule of thumb. Before Amara commits, she wants the numbers checked.

\smallskip
\begin{databox}
\textbf{Amara's numbers (everything you need)}
\begin{itemize}[leftmargin=14pt, itemsep=2pt, topsep=3pt]
\item Cost to make one box: \textbf{\$6.00}. Fixed costs for the line: \textbf{\$1{,}200 per month}.
\item Required profit margin: \textbf{40\%}. Margin is profit as a share of the \emph{price}.
\item A campus market study estimates sales of about \textbf{400 boxes per month} at a price near \$10.
\item The partner's price sheet: ``cost plus a 40\% markup, so \$6.00 + 40\% = \textbf{\$8.40} per box.''
\end{itemize}
\end{databox}

\question{1. Check the Partner's Price}
The partner insists a 40\% markup gives a 40\% margin. At \$8.40 per box, what margin does Amara actually earn? As a group, settle whether markup and margin are the same thing, and be ready to defend your answer.
\hint{Margin divides the profit by the price; markup divides it by the cost.}

\question{2. Find the Right Price}
What price gives a true 40\% margin? Verify your price by computing its margin.
\hint{At the right price, the \$6.00 cost is exactly 60\% of the price.}

\question{3. Will the Line Survive?}
For the partner's \$8.40 and for your Question 2 price: how many boxes per month does the line need just to cover the \$1{,}200 in fixed costs, and what is the monthly profit or loss at the expected 400 boxes? Which price makes the launch viable?
\hint{Each box contributes its profit toward the fixed costs; break-even is fixed costs divided by profit per box.}

\smallskip
\noindent\textit{As a group, write Amara a 2--3 sentence recommendation: launch or not, at what price, and what to tell the partner about the price sheet.}

\vfill

\noindent\rule{\textwidth}{0.4pt}
\begin{center}\textbf{ANSWER KEY}\end{center}
{\small
\textit{1)} Profit at \$8.40 is $8.40 - 6.00 = \$2.40$. Margin $= 2.40/8.40 = \boxed{28.6\%}$, well short of 40\%. Markup and margin are different: markup measures against cost, margin against price, so a 40\% markup always yields a smaller margin.

\textit{2)} Require $\text{price} - 6.00 = 0.40 \times \text{price}$, so $6.00 = 0.60 \times \text{price}$ and $\text{price} = \boxed{\$10.00}$. Check: $(10.00 - 6.00)/10.00 = 40\%$.

\textit{3)} Profit per box: \$2.40 at \$8.40; \$4.00 at \$10.00. Break-even: $1200/2.40 = \boxed{500 \text{ boxes}}$ at \$8.40 versus $1200/4.00 = \boxed{300 \text{ boxes}}$ at \$10.00. At the expected 400 boxes: \$8.40 loses $400(2.40) - 1200 = -\$240$ per month; \$10.00 earns $400(4.00) - 1200 = +\$400$ per month. Only the \$10.00 price is viable.

\textit{Verdict.} Launch at \$10.00: expected demand (400) clears the 300-box break-even and nets about \$400 a month. At the partner's \$8.40 the line loses \$240 a month at the same demand, so the price sheet must be corrected before printing menus.
}

% ================= VETTED PLAYBOOK USED (Intro Finance, from library) =================
% P1 Scope: markup vs margin + linear break-even (subtopic supplied); no calculus; formulas
%    described in plain words in hints.
% P2 Grain: two concepts, decision compares two real options (two candidate prices) against
%    a demand threshold.
% P3 Full strength: options differ structurally (a rule-of-thumb price vs a margin-derived
%    price with different break-evens), not just different numbers on identical structures;
%    break-even solved, not eyeballed.
% P4 Clean numbers: whole dollars, margin 40%, cost $6.00, break-evens 500 and 300,
%    profits -$240 and +$400 at 400 boxes.
% P5 Trap: confusing markup with margin; wrong path gives $8.40 and a money-losing launch;
%    the gap flips the decision. Framed as the partner's claim = the group discussion.
% P6 Authentic user: small-business owner pricing a product line.
% P7 Interpretation: every computed number restated as dollars per month and launch/no-launch.
%
% Instructor checklist (same_page mode, kept as comments):
%  1. Compute profit per box at $8.40 first, then divide by the PRICE for margin.
%  2. For Q2, write the margin requirement as an equation and solve for price.
%  3. Break-even = fixed costs / profit per box. Do it at both prices.
%  4. Compare both break-evens to 400 expected boxes, then compute profit at 400.
% Alignment: Intro Finance; markup vs margin; contribution and break-even analysis;
% objectives: (1) distinguish markup from margin and price to a required margin;
% (2) compute break-even from fixed costs and per-unit contribution; (3) turn both into a
% launch decision against expected demand. Inert distractors: none beyond flavor (bakery
% setting, menu printing).
% ======================================================================================
\end{document}
